\section*{Story of Michael Siems}

My first contact with a boomerang was during my childhood. The first time I threw one I had to search for it! The boomerang tended to land as far away as possible, instead of, as I had hoped, landing near me. From that moment, this mysterious object, the boomerang, simply would not leave my mind.
Most people have the same experience during their first encounter with a boomerang and for me, the subject was forgotten for a while. However, the boomerang was never far from my mind. 

One day I had a strong urge to make one myself and I asked myself, how are boomerangs made? There was no information available so once again the subject was forgotten. 
But then in April 1983, I passed a huge athletic field south of Amsterdam. My immediate thought was: I just have to throw something over it. Back home again, I explained my ambitions to a friend. As luck would have it, a few days later I was given ``Throwing and catching boomerangs'' by Günther Veit. He placed the book on the table and three hours later I’d finished making my first boomerang. That was the beginning of a new hobby, which gradually began to take shape after a slow start. In May 1983 I visited the Whitsun Tournament in Amsterdam with my first 5 boomerangs. It was there that I met Barnaby Ruhe, who was one of the talented throwers competing. It was at this championship that I was amazed to see this American Champion running and climbing trees to retrieve his boomerang! The fact that even champions had to perform these tasks improved my self-confidence so much that shortly afterward I entered the German Championship and was placed third. Since then I have participated in almost all the German and European Championships and many Tournaments abroad. I have served as Captain of the German Team four times during 9 World Championships. 

Competition performances have improved significantly during the last few years. In 1983 the record for Maximum Time Aloft (MTA) was 27 seconds (David Schummy, Great Britain). Now it is more than 2 minutes 59 seconds (Dennis Joyce, USA). Similar advances have also been recorded in other events.
\newpage
Boomerang throwing is a sport which costs next to nothing. A piece of wood, an hour to make a boomerang, a field to throw and away you go! Once you have caught on to it, you can design your own models, create new shapes and forms and decorate your boomerangs. You will meet interesting people, who out of curiousity come up and ask ``Is that a boomerang?'' In many cases they catch the ``boomerang fever''. 

Boomerang throwing is by no means a single faceted sport. You have to be able to run, throw and catch and to deal with the wind. Above all: maintain concentration — or you’ll end up with scratches and bruises! You have to be able to climb trees, crawl, dive, swim and finally: be able to search with a lot of patience, perserverance and intelligent guesswork.




